% Edited conversion of source p. 25.
\ReportHeading
  {Моделювання викривлення та втрати стійкості бурильної колони у тривимірних криволінійних свердловинах}
  {О.~М. Андрусенко}
  {Національний технічний університет України ``Київський політехнічний інститут імені Ігоря Сікорського''}
  {}

Удосконалення технологій буріння глибоких спрямованих нафтових і газових свердловин має важливе промислове значення, оскільки дає змогу підвищити продуктивність свердловин і обсяги видобутку. Основні труднощі під час проходження довгих криволінійних свердловин пов’язані з подоланням значних силових перешкод, зумовлених одночасною дією сили тяжіння, контактних сил і тертя. Ці сили залежать від технологічного режиму, жорсткості бурильної колони, звивистості стовбура та його довжини і можуть призводити до викривлення колони, втрати її стійкості та заклинювання.

Для прогнозування й запобігання таким станам потрібні спеціальні математичні моделі та методи комп’ютерного моделювання. Водночас механічні явища у довгих криволінійних свердловинах є складними, тому в інженерній практиці часто застосовують істотно спрощені підходи.

У роботі розглянуто моделювання режимів буріння глибоких криволінійних свердловин із заданими недосконалими геометричними траєкторіями осьових ліній. На основі теорії криволінійних гнучких пружних стрижнів, диференціальної геометрії та чисельного аналізу розроблено тривимірну жорсткісну модель тяги та крутного моменту бурильної колони, а також програмне забезпечення для моделювання режимів спуску, підйому і буріння.

Комп’ютерні розрахунки показують, що контактні та фрикційні сили можна визначати й регулювати так, щоб забезпечувати заздалегідь спроєктовані безаварійні режими роботи. Розроблені моделі та програмні засоби можуть бути використані для прогнозування аварійних ситуацій і запобігання їм на стадіях проєктування та реалізації процесу буріння.

\begin{thebibliography}{9}
\bibitem{drill:gulyayev2019}
V.~Gulyayev, S.~Glazunov, O.~Glushakova, E.~Vashchilina, L.~Shevchuk, N.~Shlyun, and E.~Andrusenko,
\emph{Modelling Emergency Situations in the Drilling of Deep Boreholes}.
Cambridge: Cambridge Scholars Publishing, 2019, 544~p.

\bibitem{drill:gulyayev2017}
V.~I. Gulyayev, E.~N. Andrusenko, and S.~N. Glazunov,
``Computer simulation of resistance-force mitigation through curvature bridging in extended boreholes,''
\emph{Journal of Petroleum Science and Engineering}, vol.~156, pp.~594--604, 2017.
\end{thebibliography}
